% Adapted from Jake's Resume (MIT License)
\documentclass[letterpaper,10pt]{article}

\usepackage[T1]{fontenc}
\usepackage[utf8]{inputenc}
\usepackage{lmodern}
\usepackage[empty]{fullpage}
\usepackage{titlesec}
\usepackage{enumitem}
\usepackage[hidelinks]{hyperref}
\usepackage{fancyhdr}
\usepackage{tabularx}
\usepackage{xcolor}
\usepackage{microtype}
\input{glyphtounicode}
\pdfgentounicode=1

\pagestyle{fancy}
\fancyhf{}
\renewcommand{\headrulewidth}{0pt}
\renewcommand{\footrulewidth}{0pt}

\addtolength{\oddsidemargin}{-0.62in}
\addtolength{\evensidemargin}{-0.62in}
\addtolength{\textwidth}{1.24in}
\addtolength{\topmargin}{-0.66in}
\addtolength{\textheight}{1.30in}

\urlstyle{same}
\raggedbottom
\raggedright
\setlength{\tabcolsep}{0in}
\setlength{\parindent}{0pt}

\titleformat{\section}{\vspace{-3pt}\scshape\raggedright\large}{}{0em}{}[\color{black}\titlerule \vspace{-5pt}]

\newcommand{\resumeItem}[1]{\item\small{#1\vspace{-1pt}}}
\newcommand{\resumeSubheading}[4]{
  \vspace{0pt}\item
  \begin{tabular*}{0.98\textwidth}[t]{l@{\extracolsep{\fill}}r}
    \textbf{#1} & #2 \\
    \textit{\small #3} & \textit{\small #4} \\
  \end{tabular*}\vspace{-6pt}
}
\newcommand{\resumeSubSubheading}[2]{
  \item
  \begin{tabular*}{0.98\textwidth}{l@{\extracolsep{\fill}}r}
    \textit{\small #1} & \textit{\small #2} \\
  \end{tabular*}\vspace{-6pt}
}
\newcommand{\resumeProjectHeading}[2]{
  \item
  \begin{tabular*}{0.98\textwidth}{l@{\extracolsep{\fill}}r}
    \small #1 & \small #2 \\
  \end{tabular*}\vspace{-6pt}
}
\renewcommand\labelitemii{$\vcenter{\hbox{\tiny$\bullet$}}$}
\newcommand{\resumeSubHeadingListStart}{\begin{itemize}[leftmargin=0.14in,label={}]}
\newcommand{\resumeSubHeadingListEnd}{\end{itemize}}
\newcommand{\resumeItemListStart}{\begin{itemize}[leftmargin=0.22in,itemsep=1pt,topsep=2pt]}
\newcommand{\resumeItemListEnd}{\end{itemize}\vspace{-5pt}}

\hypersetup{
  pdftitle={Arseniy Mironov - Resume},
  pdfauthor={Arseniy Mironov},
  pdfsubject={Machine Learning and LLM Inference Resume}
}

\begin{document}

\begin{center}
  \textbf{\Huge \scshape Arseniy Mironov} \\ \vspace{2pt}
  \small Nizhny Novgorod, Russia $|$ +7 (917) 618-94-15 $|$
  \href{mailto:arseniy.mironov.dev@gmail.com}{\underline{arseniy.mironov.dev@gmail.com}} \\
  \href{https://github.com/Napkin-AI}{\underline{github.com/Napkin-AI}} $|$
  \href{https://t.me/xNapkin}{\underline{t.me/xNapkin}}
\end{center}

\section{Education}
\resumeSubHeadingListStart
  \resumeSubheading
    {Lobachevsky State University of Nizhny Novgorod (UNN)}{Nizhny Novgorod, Russia}
    {M.Sc. in Applied Mathematics and Informatics}{2026 -- Present}
  \resumeSubSubheading
    {B.Sc. in Fundamental Informatics and Information Technologies}{2022 --- 2026}
  \resumeSubheading
    {Yandex School of Data Analysis}{ }
    {Data Science Track}{2nd year}
\resumeSubHeadingListEnd

\section{Experience}
\resumeSubHeadingListStart
  \resumeSubheading
    {Engineer}{Jan. 2026 -- Present}
    {Huawei}{Nizhny Novgorod, Russia}
  \resumeItemListStart
    \resumeItem{Develop the SGLang inference stack for Ascend NPUs, focusing on LLM and diffusion-model inference, attention backends, and performance optimization.}
    \resumeItem{Integrated sparse attention backends into SGLang Diffusion for Ascend NPU, including backend configuration, accuracy validation, documentation, and performance profiling; merged in \href{https://github.com/sgl-project/sglang/pull/23482}{\underline{Pull Request \#23482}}.}
    \resumeItem{Investigated Qwen3 and diffusion-LLM prefill acceleration and inference bottlenecks on Ascend NPUs.}
  \resumeItemListEnd

  \resumeSubheading
    {Engineer Intern}{Jul. 2025 -- Aug. 2025}
    {Huawei}{Nizhny Novgorod, Russia}
  \resumeItemListStart
    \resumeItem{Researched disaggregated execution of prefill and decode stages in vLLM on hardware accelerators.}
    \resumeItem{Adapted the KV-cache transfer engine to the target accelerator architecture and benchmarked LLM performance under disaggregated serving.}
  \resumeItemListEnd
\resumeSubHeadingListEnd

\section{Open Source \& Projects}
\resumeSubHeadingListStart
  \resumeProjectHeading
    {\textbf{OpenCV - QR and ArUco Detection Optimization} $|$ \emph{C++}}{}
  \resumeItemListStart
    \resumeItem{Implemented \texttt{cv::approxPolyN}, a greedy convex-polygon approximation algorithm with a configurable target side count and area-error threshold for contour processing.}
    \resumeItem{The feature \href{https://github.com/opencv/opencv/pull/25607}{\underline{Pull Request \#25607}} was merged into OpenCV and shipped in OpenCV 4.11.0.}
  \resumeItemListEnd
\resumeSubHeadingListEnd


\section{Publications}
\resumeSubHeadingListStart
  \resumeProjectHeading
    {\textbf{A. I. Mironov}, N. A. Sokolov, A. A. Serebryakova}{2025}
  \resumeItemListStart
    \resumeItem{``The Effect of the Volume of Synthetically Generated Images in the Dataset for the Task of Axon Segmentation in Brain Electron Microscopy.'' \textit{GraphiCon 2025}. \href{https://doi.org/10.25686/978-5-8158-2474-4-2025-828-834}{\underline{DOI}}}
  \resumeItemListEnd
\resumeSubHeadingListEnd

\section{Competitive Programming}
\resumeSubHeadingListStart
  \resumeProjectHeading{\textbf{ICPC 2025--2026}}{}
  \resumeItemListStart
    \resumeItem{Nizhny Novgorod qualifier: 2nd-degree diploma; Saratov quarterfinal: 3rd-degree diploma; ICPC semifinal, NERC final in St. Petersburg: 3rd-degree diploma.}
  \resumeItemListEnd
  \resumeProjectHeading{\textbf{ICPC 2024--2025}}{}
  \resumeItemListStart
    \resumeItem{Nizhny Novgorod qualifier: 2nd-degree diploma; Saratov quarterfinal: 59th place.}
  \resumeItemListEnd
\resumeSubHeadingListEnd

\section{Technical Skills}
\begin{itemize}[leftmargin=0.14in,label={}]
  \small{\item{
    \textbf{Languages:} Python, C++, SQL, Triton \\
    \textbf{ML / Inference:} PyTorch, SGLang, vLLM, OpenCV, NumPy, pandas, scikit-learn, Matplotlib \\
    \textbf{Systems / Tools:} Linux, Git, Jupyter, LaTeX \\
    \textbf{English:} B1
  }}
\end{itemize}

\end{document}
